\chapter{Overview}
\label{chap:overview}

A quadrotor carries a single camera on a one-axis pitch hinge and must keep an
airborne target in view in \emph{every} camera frame. Its only measurements of
the target are a bearing and an apparent size (a bounding box), from which range
follows because the target's true dimensions are known.

The camera is the binding constraint. A quadrotor can only accelerate by tilting
its thrust axis, and the camera is attached to the body, so every manoeuvre the
drone makes in order to keep up also turns the camera. On top of that, the field of
view is a cone of fixed half-angle, the pixel budget caps the range, and every
rotation (tilt, yaw, hinge) is rate-limited.

\paragraph{Goal (P3).} One explicit \emph{master inequality} in the parameters
\[
  (\text{FOV},\ \text{resolution},\ \text{hinge range and rate},\ \text{tilt and
  yaw rates},\ \text{thrust},\ \text{drag},\ \delta;\ \text{target speed band},
  \ \text{turn radius},\ \text{size})
\]
that decides whether tracking can be maintained. It should be sufficient as
stated and necessary up to an explicitly bounded gap that vanishes as the control
period shrinks and the rate limits grow (a goal: the current sufficient side needs a fixed-time-window filter for this). From it we read off design rules such as
``to hold a standoff of $d$ metres against this target you need a half-FOV of at
least \dots'' and a map of regimes (fixed camera versus hinge, fast versus slow
drone).

\paragraph{Route.}
\begin{enumerate}
  \item \textbf{Model} (Chapter~\ref{chap:model}): discrete-time dynamics of target
        and drone, the camera geometry, visibility, and strategies.
  \item \textbf{Hard limits} (Chapter~\ref{chap:speed}): if the target is faster
        than the drone, no strategy works. \emph{Proved.}
  \item \textbf{Pointing reduction} (Chapter~\ref{chap:pointing}, \emph{proved}): with yaw free, ``can the camera see direction
        $u$'' collapses to one scalar, the elevation of $u$ above the rotor plane.
  \item \textbf{Replication and formation hold}: sufficient conditions under which
        the drone copies the target's motion at a fixed offset, with explicit lag
        terms.
  \item \textbf{Converse bounds}: target manoeuvres that defeat every strategy
        outside the region of item~4.
  \item \textbf{Master inequality} assembling items 2--5.
\end{enumerate}

\paragraph{Reference scenario.} Target: a fixed-wing aircraft of length
$3.5$\,m, wingspan $2.5$\,m, height $0.5$\,m, cruise $40$\,m/s, top speed
$51.4$\,m/s, minimum turn radius $350$\,m (so lateral acceleration
$v^2/R \approx 4.6$--$7.5$\,m/s$^2$, turn rate $\le 0.15$\,rad/s), assumed stall
speed $20$\,m/s. Detection size $s = 2.5$\,m (the smallest longer side of its
bounding box over all viewing directions). With $1920$ pixels across a $90^\circ$
field of view and an $8$-pixel threshold, the detection range is about $382$\,m.
Every drone and camera quantity stays a free parameter; numbers are only for
illustration.
